\section{Spherical monads}\label{appendix:sphmonad}

\subsection{Strengthening Christ's theorem}\label{subsection:christ}
We recall the following theorem of Christ.

\begin{thm}[Christ, \cite{christSphericalMonadicAdjunctions2023}*{Theorem 4.1}]\label{thm:christ}
    Let $\Phi$ be a stable $\infty$-category and $M$ a monad with unit $u:\id\to M$ and let $T:=\Cone(u)$. The following are equivalent.
    \begin{enumerate}
        \item \label{item:1} The endofunctor $T$ is an equivalence, and the unit $u$ satisfies $T u\simeq uT$.
        \item \label{item:2} The monadic adjunction $S:\Phi\tofrom \Mod_\Phi(M):R$ is spherical.
    \end{enumerate}
\end{thm}

Let us split statement (\ref{item:1}) into two parts.
\begin{enumerate}[label=(\roman*)]
    \item \label{item:1a} The endofunctor $T$ is an equivalence.
    \item \label{item:1b} The unit $u$ satisfies $T u\simeq uT$.
\end{enumerate}

Theorem~\ref{thm:christstrengthening1} strengthens this by showing that statement \ref{item:1b} in fact always holds. We begin by recalling the meaning of $Tu\simeq uT$. As written here, this is ambiguous as these have different targets, $TM$ and $MT$ respectively. To make sense of this, we recall a version of \cite{christSphericalMonadicAdjunctions2023}*{Lemma 2.2}, stated in terms of monads instead of adjunctions.

\begin{lem}\label{lem:commutes}
    Let $M$ be a monad with unit $u$ and let $T:=\Cone(u:\id\to M)$. Then, there is a canonical identification $TM\simeq MT$.
\end{lem}

\begin{proof}
    Consider the naturally commuting square
    \[\begin{tikzcd}
        M & MM \\
        M & M.
        \arrow["uM", from=1-1, to=1-2]
        \arrow["{1_M}", from=1-1, to=2-1]
        \arrow["m", from=1-2, to=2-2]
        \arrow["{1_M}", from=2-1, to=2-2]
    \end{tikzcd}\]
    We may compute the total complex in two ways. Taking the fiber vertically first and then cone horizontally yields $\coCone(m)$. Taking cone horizontally and then fiber vertically yields $TM$. So we have the identification,
    \[\coCone(m)\simeq TM.\]
    By symmetry, we also have
    \[\coCone(m)\simeq MT.\]
    Composing these two provides the desired identification, $TM\simeq MT$.
\end{proof}

We may now state the key assertion.

\begin{prop}\label{prop:main}
    Let $\Phi$ be a stable $\infty$-category and $M$ a monad with unit $u:\id\to M$ and let $T:=\Cone(u)$. Then, the unit $u$ satisfies $Tu\simeq uT$ under the equivalence $TM\simeq MT$ of Lemma~\ref{lem:commutes}.
\end{prop}

\begin{proof}
    The idea of the proof is the same as, and will be clearly compatible with the proof of Lemma~\ref{lem:commutes}. Consider the following commutative cube which we call $X$.
    % https://q.uiver.app/#q=WzAsOSxbMSwxLCJNIl0sWzMsMSwiTU0iXSxbMSwzLCJNIl0sWzMsMywiTSJdLFszLDJdLFswLDAsIjEiXSxbMCwyLCJNIl0sWzIsMCwiTSJdLFsyLDIsIk0iXSxbMCwxLCJ1TSJdLFswLDJdLFsyLDNdLFsxLDMsIm0iXSxbNSwwLCJ1Il0sWzYsMl0sWzcsMSwiTXUiXSxbOCwzXSxbNSw3LCJ1Il0sWzUsNiwidSJdLFs2LDhdLFs3LDhdXQ==
    \[\begin{tikzcd}
        1 && M & \\
        & M && MM \\
        M && M & {} \\
        & M && M
        \arrow["u", from=1-1, to=1-3]
        \arrow["u", from=1-1, to=2-2]
        \arrow["u", from=1-1, to=3-1]
        \arrow["Mu", from=1-3, to=2-4]
        \arrow["{{ }}"{description}, from=1-3, to=3-3]
        \arrow["uM", from=2-2, to=2-4]
        \arrow[from=2-2, to=4-2]
        \arrow["m", from=2-4, to=4-4]
        \arrow[from=3-1, to=3-3]
        \arrow[from=3-1, to=4-2]
        \arrow[from=3-3, to=4-4]
        \arrow[from=4-2, to=4-4]
    \end{tikzcd}\]
    For each direction $*\in\{\hor,\ver,\diag\}$, we denote by $\Cone_*(-)$ the result of applying $\Cone$ in the direction indicated by $*$. We write $\coCone_*(-)$ similarly.

    We have comparisons,
    \begin{equation}\label{eqn:comparison1}
        \coCone_{\ver}(\Cone_{\hor}(X))\simeq\Cone_{\hor}(\coCone_{\ver}(X)),
    \end{equation}
    and
    \begin{equation}\label{eqn:comparison2}
        \coCone_{\ver}(\Cone_{\dia}(X))\simeq\Cone_{\dia}(\coCone_{\ver}(X)).
    \end{equation}
    It is easy to calculate that
    \[\coCone_{\ver}(\Cone_{\hor}(X))\simeq\begin{tikzcd}
        T & \\
        & TM.
        \arrow["Tu", from=1-1, to=2-2]
    \end{tikzcd}\]
    and similarly that
    \[\coCone_{\ver}(\Cone_{\dia}(X))\simeq
    % https://q.uiver.app/#q=WzAsMixbMCwwLCJUIl0sWzIsMCwiTVQiXSxbMCwxLCJ1VCJdXQ==
    \begin{tikzcd}
        T && MT.
        \arrow["uT", from=1-1, to=1-3]
    \end{tikzcd}\]
    We also calculate that
    % https://q.uiver.app/#q=WzAsNCxbMCwwLCJUWy0xXSJdLFsxLDEsIjAiXSxbMywxLCJcXEZpYihtKSJdLFsyLDAsIjAiXSxbMCwxXSxbMywyXSxbMSwyXSxbMCwzXV0=
    \[\coCone_{\ver}(X)\simeq\begin{tikzcd}
        {T[-1]} && 0 & \\
        & 0 && {\coCone(m)}.
        \arrow[from=1-1, to=1-3]
        \arrow[from=1-1, to=2-2]
        \arrow[from=1-3, to=2-4]
        \arrow[from=2-2, to=2-4]
    \end{tikzcd}\]
    So we have the comparison,
    \begin{equation}\label{eqn:comparison3}
        \Cone_{\hor}(\coCone_{\ver}(X))\simeq\Cone_{\dia}(\coCone_{\ver}(X)).
    \end{equation}
    Composing comparisons (\ref{eqn:comparison1}),(\ref{eqn:comparison2}) and (\ref{eqn:comparison3}), and taking care of signs, we deduce the desired result.
\end{proof}

\begin{proof}[Proof of Theorem~\ref{thm:christstrengthening1}]
    This follows from the original version, Theorem~\ref{thm:christ}, together with Proposition~\ref{prop:main}.
\end{proof}

\subsection{A description of objects of \texorpdfstring{$\sfSphMnd(n)$}{SphMnd(n)}}\label{subsection:sphmonad}
We recall that currently, $\sfSphMnd(n)$ is defined as
\[\sfSphMnd(n) := \sfAut(n)\times_{\sfAut}\sfSphMnd.\]
In this subsection, we provide a direct characterization of abstract prelocalized perverse schobers, analogously to and compatibly with the equivalence $\sfAdj(n)\simeq\sfBdj(n)$ of Proposition~\ref{prop:bdjisadj}.
%As a result, we construct a projection,
%\[\sftwoP(\bbC,R)\to \overline{\sftwoP}^{\opname{pre}}(\bbC,R)\]
%left adjoint to the natural inclusion.

The following corollary of Proposition~\ref{prop:conservative} says that we are allowed to forget about $\Psi$.

\begin{cor}\label{cor:sphmonad}
    A prelocalized schober is equivalent to the data of categories $\Phi_i$ together with a monad $M$ on $\oplus_i\Phi_i$ such that for each $j$, the composition
    \[\Phi_j \to\oplus_i\Phi_i \to \Mod_M(\oplus_i\Phi_i)\]
    is spherical.
\end{cor}

\begin{proof}
    We apply Barr-Beck to the adjunction,
    \[\oplus_i\Phi_i \tofrom \Psi.\]
    The right adjoint is conservative by Proposition~\ref{prop:conservative} and spherical functors automatically preserve colimits. Moreover, all categories are assumed cocomplete. We deduce that there is an identification
    \[\Psi\isomto\Mod_M(\oplus_i\Phi_i)\]
    thus we can reconstruct the perverse schober from just the data of the monad.
\end{proof}

\begin{defn}
    A monad $M=\oplus_{i,j}M_{ij}$ acting on $\oplus_i\Phi_i$ is said to be a spherical matrix monad if the condition in Corollary~\ref{cor:sphmonad} is satisfied.
\end{defn}

It is interesting to characterize which monads on $\oplus_i\Phi_i$ are spherical matrix monads without reference to the Eilenberg--Moore adjunction, as we studied for $n=1$ in the previous section. The following is a generalization of Theorem~\ref{thm:christstrengthening1}. We may also understand this as a categorification of the quiver description in \cite{gelfandPerverseSheavesQuivers1996}.

\begin{prop}\label{prop:sphmonad}
    Consider categories $\Phi_i$, together with a monad $M=\oplus_{i,j}M_{ij}$ acting on $\oplus_i\Phi_i$. This monad is a spherical matrix monad if and only if the following two conditions are satisfied.
    \begin{enumerate}
        \item \label{enum:sphmonad1} For each $i$, $T_i:=\coCone(\id_{\Phi_i}\to M_{ii})$ is an equivalence of categories.
        \item \label{enum:sphmonad2} For each $i,j$, the map $M_{ji}\to M_{ij}^RT_i[1]$ is an equivalence of functors, where this map is obtained via adjunction from the composition $M_{ij}M_{ji}\to M_{ii}\to T_i[1]$.
    \end{enumerate}
    In fact, we may replace Condition~(\ref{enum:sphmonad2}) with the following weaker condition.
    \begin{enumerate}[resume]
        \item \label{enum:sphmonad3} For each $i\neq j$, the map $M_{ji}\to M_{ij}^RT_i[1]$ is an equivalence of functors, where this map is obtained via adjunction from the composition $M_{ij}M_{ji}\to M_{ii}\to T_i[1]$.
    \end{enumerate}
\end{prop}

\begin{proof}
    The stronger statement follows from the weaker statement. Indeed, by Theorem~\ref{thm:christstrengthening1}, Condition~(\ref{enum:sphmonad1}) implies that for each $i$, $M_{ii}$ is a spherical monad. This then implies Condition~(\ref{enum:sphmonad2}) whenever $j=i$.
    
    We prove the weaker statement. We begin by noting that the adjunctions
    \[S_i:\Phi_i\tofrom\opname{Mod}_{\oplus_i\Phi}(M_{ij}):R_i\] 
    recover the monad via $R_iS_j\isomto M_{ij}$. In particular, $T_i \simeq \coCone(\id\to R_iS_i)$.

    Suppose that $M$ is a spherical matrix monad, or equivalently that each $S_i$ is spherical. Then by the above, Condition~(\ref{enum:sphmonad1}) is clear. We prove Condition~(\ref{enum:sphmonad2}). Recall that $R_i$ admits a right adjoint $R_i^R$ and the natural map
    \[S_i\to R^R_iT_i[1]\]
    is an equivalence. Composing with $R_j$ on the left, we see that $M_{ji}\to M_{ij}^RT_i[1]$ is an equivalence.
    
    We must be careful here and verify that this is identified with the natural transformation coming from the composition $M_{ij}M_{ji}\to M_{ii}\to T_i[1]$ by adjunction. We observe that $M_{ii}\to T_i[1]$ can be post-composed after the fact, so it suffices to compare the natural transformations
    \begin{itemize}
        \item $R_j(S_i\to (R_i^RR_i)S_i)$, and
        \item $M_{ji}\to M_{ij}^RM_{ii}$.
    \end{itemize}
    This can be seen by comparing the following diagrams representing these natural transformations, where the functors are composed horizontally as in writing order, and natural transformations go downwards.
    %Comments: The \path command is there to define all of the points I will need to use. Then, \draw draws the curves, and \fill draws the 2-cells. There is also a way to add dots, but I don't need them here.

\tikzstyle{catc}=[red!30]
\tikzstyle{catd}=[blue!30]

\[
\begin{tikzpicture}[scale=0.75]
\path coordinate[] (epsilon)
+(-1.5,-1.5) coordinate[label=below:$R_i^R$] (tl)
+(1.5,-1.5) coordinate[label=below:$R_i$] (tr)
+(-2.5,1.5) coordinate[label=above:$R_j$] (ul)
+(2.5,1.5) coordinate[label=above:$S_i$] (ur)
+(-2.5,-1.5) coordinate[label=below:$R_j$] (dl)
+(2.5,-1.5) coordinate[label=below:$S_i$] (dr);
\draw (tl) to[out=90, in=180] (epsilon.center) to[out=0, in=90] (tr);
\draw (ul) to (dl);
\draw (ur) to (dr);
\begin{scope}[on background layer]
\fill[catd] +(-3.5,-1.5) rectangle +(3.5,1.5);
\fill[catc] (tl) to[out=90, in=180] (epsilon.center) to[out=0, in=90] (tr) --
cycle;
\fill[catc] +(-3.5,-1.5) rectangle +(-2.5,1.5);
\fill[catc] +(2.5,-1.5) rectangle +(3.5,1.5);
\end{scope}
\end{tikzpicture}
\quad
\begin{tikzpicture}[scale=0.75]
\path coordinate[] (tr)
++(0,-3) coordinate[] (br)
++(-1,0) coordinate[label=below:$S_i$] (brl)
++(0,3) coordinate[label=above:$S_i$] (trl)
++(-1,0) coordinate[label=above:$R_j$] (a)
++(0,-2) coordinate[] (b)
++(-0.5,-0.5) coordinate[] (c)
++(-0.5,0.5) coordinate[] (d)
++(-1.5,1.5) coordinate[] (e)
++(-1.5,-1.5) coordinate[] (f)
++(0,-1) coordinate[label=below:$R_j$] (g)
++(-1,0) coordinate[] (bl)
++(0,3) coordinate[] (tl)
++(2,-3) coordinate[label=below:$R_i^R$] (h)
++(0,1) coordinate[] (i)
++(0.5,0.5) coordinate[] (j)
++(0.5,-0.5) coordinate[] (k)
++(0,-1) coordinate[label=below:$R_i$] (l);
\draw (trl) -- (brl);
\draw (a) -- (b) to [out=-90, in=0] (c.center) to [out=180, in=-90] (d) to [out=90, in=0] (e.center) to [out=180, in=90] (f) -- (g);
\draw (h) -- (i) to [out=90, in=180] (j.center) to [out=0, in=90] (k) -- (l);
\begin{scope}[on background layer]
\fill[catd] (bl) rectangle (tr);
\fill[catc] (h) -- (i) to [out=90, in=180] (j.center) to[out=0, in=90] (k) -- (l) -- cycle;
\fill[catc] (a) -- (b) to [out=-90, in=0] (c.center) to [out=180, in=-90] (d) to [out=90, in=0] (e.center) to [out=180, in=90] (f) -- (g) -- (bl) -- (tl) -- cycle;
\fill[catc] (brl) rectangle (tr);
\end{scope}
\end{tikzpicture}
\]
    Using this diagram as a mnemonic, we may produce the desired identification.

    Let us now prove that if $M$ satisfies Condition~(\ref{enum:sphmonad1}) and Condition~(\ref{enum:sphmonad2}), then for each $i$, $S_i$ is spherical. We will apply the 2-out-of-4 property, \cite{christSphericalMonadicAdjunctions2023}*{Theorem 2.15}. We must first show that $R_i$ admits a right adjoint. It is clear that $M_{ii}$ admits a right adjoint $M_{ii}^R$, since $M$ is an extension of two autoequivalences which each automatically admit right adjoints. Observing that $\opname{Mod}_{\Phi}(M_{ii})\isomto\opname{coMod}_{\Phi}(M_{ii}^R)$ we get the right adjoint from the comonadic adjunction,
    \[R_i: \opname{coMod}_{\Phi}(M_{ii}^R)\tofrom \Phi: R_i^R.\]
    
    We can now apply the 2-out-of-4 property. It is immediate that the cotwist functor of $S_i\adjto R_i$ is invertible, and so it suffices to verify that the natural map $S_i\to R_i^RT_i[1]$ is an equivalence. This becomes an equivalence after we post-compose with $\oplus_jR_j$, by Condition~(\ref{enum:sphmonad2}) and the comparison made above. Since $\oplus_jR_j$ is conservative, it follows that $S_i\to R_i^RT_i[1]$ is indeed an equivalence.
\end{proof}

We may also construct the projection functor directly.

\begin{cor}
    Given an object $(\Phi_i,\Psi,S_i,R_i)$ of $\sfSph(n)$, the monad $M=\oplus_{i,j}R_iS_j$ is a spherical matrix monad. In particular, there is a projection functor 
    \[\sfSph(n)\to \sfSphMnd(n).\]
    This agrees with the projection functor constructed in \S\ref{subsection:quotientismonad}.
\end{cor}

\begin{proof}
    The two conditions of Proposition~\ref{prop:sphmonad} follow by the same argument as the proof of the ``only if'' direction, noting that this direction didn't use the conservativity of $\oplus_i R_i$.
\end{proof}

\begin{cor}\label{prop:decategorification2}
    Recall notation from Proposition~\ref{prop:gmvquiver}. Given an abstract prelocalized perverse schober given as a spherical matrix monad $\oplus_{i,j} M_{ij}$ acting on $\oplus_i\Phi_i$, we may form an object of $\opname{Q}(n)$ by taking $M_i := K_0(\Phi_i)$ and $m_{ij} = K_0(M_{ij})$.
\end{cor}
